\documentclass[11pt,compress,t,notes=noshow, xcolor=table]{beamer}
\usepackage[]{graphicx}\usepackage[]{color}
% maxwidth is the original width if it is less than linewidth
% otherwise use linewidth (to make sure the graphics do not exceed the margin)
\makeatletter
\def\maxwidth{ %
  \ifdim\Gin@nat@width>\linewidth
    \linewidth
  \else
    \Gin@nat@width
  \fi
}
\makeatother

\newcommand{\citebutton}[2]{%
\beamergotobutton{\href{#2}{#1}}%
}

\newcommand{\blu}[1]{\textcolor{blue}{#1}}
\newcommand{\org}[1]{\textcolor{orange}{#1}}
\newcommand{\ques}{\textbf{\textcolor{red}{Question:  }}}
\newcommand{\questionssofar}{\begin{frame}\frametitle{Any questions?}\end{frame}}

\newcommand\warning{%
 \makebox[1.4em][c]{%
 \makebox[0pt][c]{\raisebox{.1em}{\scriptsize!}}%
 \makebox[0pt][c]{\color{red}\normalsize$\bigtriangleup$}}}%

\definecolor{fgcolor}{rgb}{0.345, 0.345, 0.345}
\newcommand{\hlnum}[1]{\textcolor[rgb]{0.686,0.059,0.569}{#1}}%
\newcommand{\hlstr}[1]{\textcolor[rgb]{0.192,0.494,0.8}{#1}}%
\newcommand{\hlcom}[1]{\textcolor[rgb]{0.678,0.584,0.686}{\textit{#1}}}%
\newcommand{\hlopt}[1]{\textcolor[rgb]{0,0,0}{#1}}%
\newcommand{\hlstd}[1]{\textcolor[rgb]{0.345,0.345,0.345}{#1}}%
\newcommand{\hlkwa}[1]{\textcolor[rgb]{0.161,0.373,0.58}{\textbf{#1}}}%
\newcommand{\hlkwb}[1]{\textcolor[rgb]{0.69,0.353,0.396}{#1}}%
\newcommand{\hlkwc}[1]{\textcolor[rgb]{0.333,0.667,0.333}{#1}}%
\newcommand{\hlkwd}[1]{\textcolor[rgb]{0.737,0.353,0.396}{\textbf{#1}}}%
\let\hlipl\hlkwb

\usepackage{framed}
\makeatletter
\newenvironment{kframe}{%
 \def\at@end@of@kframe{}%
 \ifinner\ifhmode%
  \def\at@end@of@kframe{\end{minipage}}%
  \begin{minipage}{\columnwidth}%
 \fi\fi%
 \def\FrameCommand##1{\hskip\@totalleftmargin \hskip-\fboxsep
 \colorbox{shadecolor}{##1}\hskip-\fboxsep
     % There is no \\@totalrightmargin, so:
     \hskip-\linewidth \hskip-\@totalleftmargin \hskip\columnwidth}%
 \MakeFramed {\advance\hsize-\width
   \@totalleftmargin\z@ \linewidth\hsize
   \@setminipage}}%
 {\par\unskip\endMakeFramed%
 \at@end@of@kframe}
\makeatother

\definecolor{shadecolor}{rgb}{.97, .97, .97}
\definecolor{messagecolor}{rgb}{0, 0, 0}
\definecolor{warningcolor}{rgb}{1, 0, 1}
\definecolor{errorcolor}{rgb}{1, 0, 0}
\newenvironment{knitrout}{}{} % an empty environment to be redefined in TeX

\usepackage{alltt}
\newcommand{\SweaveOpts}[1]{}  % do not interfere with LaTeX
\newcommand{\SweaveInput}[1]{} % because they are not real TeX commands
\newcommand{\Sexpr}[1]{}       % will only be parsed by R
\newcommand{\xmark}{\ding{55}}%


\usepackage[english]{babel}
\usepackage[utf8]{inputenc}

\usepackage{dsfont}
\usepackage{verbatim}
\usepackage{amsmath}
\usepackage{amsfonts}
\usepackage{amssymb}
\usepackage{bm}
\usepackage{csquotes}
\usepackage{multirow}
\usepackage{longtable}
\usepackage{booktabs}
\usepackage{enumerate}
\usepackage[absolute,overlay]{textpos}
\usepackage{psfrag}
\usepackage{algorithm}
\usepackage{algpseudocode}
\usepackage{eqnarray}
\usepackage{arydshln}
\usepackage{tabularx}
\usepackage{placeins}
\usepackage{tikz}
\usepackage{setspace}
\usepackage{colortbl}
\usepackage{mathtools}
\usepackage{wrapfig}
\usepackage{bm}
\usepackage{amsmath}
\usepackage{pifont}

\usetikzlibrary{shapes.multipart,shapes,arrows,automata,positioning,calc,chains,trees, shadows}
\tikzset{
  %Define standard arrow tip
  >=stealth',
  %Define style for boxes
  punkt/.style={
    rectangle,
    rounded corners,
    draw=black, very thick,
    text width=6.5em,
    minimum height=2em,
    text centered},
  % Define arrow style
  pil/.style={
    ->,
    thick,
    shorten <=2pt,
    shorten >=2pt,}
}

\tikzstyle{vec}=[draw, rectangle, fill = white, minimum width=5mm, minimum height=1cm, inner sep = 2pt]

\usepackage{subfig}

% Defines macros and environments
\usepackage{../../style/lmu-lecture}


\let\code=\texttt
\let\proglang=\textsf

\setkeys{Gin}{width=0.9\textwidth}

\setbeamertemplate{frametitle}{\expandafter\uppercase\expandafter\insertframetitle}

\usepackage{bbm}
% basic latex stuff
\newcommand{\pkg}[1]{{\fontseries{b}\selectfont #1}} %fontstyle for R packages
\newcommand{\lz}{\vspace{0.5cm}} %vertical space
\newcommand{\dlz}{\vspace{1cm}} %double vertical space
\newcommand{\oneliner}[1] % Oneliner for important statements
{\begin{block}{}\begin{center}\begin{Large}#1\end{Large}\end{center}\end{block}}


%new environments
\newenvironment{vbframe}  %frame with breaks and verbatim
{
 \begin{frame}[containsverbatim,allowframebreaks]
}
{
\end{frame}
}

\newenvironment{vframe}  %frame with verbatim without breaks (to avoid numbering one slided frames)
{
 \begin{frame}[containsverbatim]
}
{
\end{frame}
}

\newenvironment{blocki}[1]   % itemize block
{
 \begin{block}{#1}\begin{itemize}
}
{
\end{itemize}\end{block}
}

\newenvironment{fragileframe}[2]{  %fragile frame with framebreaks
\begin{frame}[allowframebreaks, fragile, environment = fragileframe]
\frametitle{#1}
#2}
{\end{frame}}


\newcommand{\myframe}[2]{  %short for frame with framebreaks
\begin{frame}[allowframebreaks]
\frametitle{#1}
#2
\end{frame}}

\newcommand{\remark}[1]{
  \textbf{Remark:} #1
}


\newenvironment{deleteframe}
{
\begingroup
\usebackgroundtemplate{\includegraphics[width=\paperwidth,height=\paperheight]{../style/color/red.png}}
 \begin{frame}
}
{
\end{frame}
\endgroup
}
\newenvironment{simplifyframe}
{
\begingroup
\usebackgroundtemplate{\includegraphics[width=\paperwidth,height=\paperheight]{../style/color/yellow.png}}
 \begin{frame}
}
{
\end{frame}
\endgroup
}\newenvironment{draftframe}
{
\begingroup
\usebackgroundtemplate{\includegraphics[width=\paperwidth,height=\paperheight]{../style/color/green.jpg}}
 \begin{frame}
}
{
\end{frame}
\endgroup
}
% https://tex.stackexchange.com/a/261480: textcolor that works in mathmode
\makeatletter
\renewcommand*{\@textcolor}[3]{%
  \protect\leavevmode
  \begingroup
    \color#1{#2}#3%
  \endgroup
}
\makeatother





\input{../../latex-math/basic-math.tex}
\input{../../latex-math/basic-ml.tex}

\newcommand{\titlefigure}{figure/transformer_sq.png}
\newcommand{\learninggoals}{
\item Understand inner workings of BPE
\item Being able to compare BPE to other tokenization approaches}

\title{Transformer}
% \author{}
\institute{\href{https://slds-lmu.github.io/lecture_dl4nlp/}{slds-lmu.github.io/lecture\_dl4nlp}}
\date{}

\begin{document}
\lecturechapter{BytePair encoding (BPE)}
\lecture{Deep Learning for NLP}

% ------------------------------------------------------------------------------

\begin{vbframe}{BytePair encoding (BPE)}

\vfill

\textbf{Data compression algorithm \href{https://www.derczynski.com/papers/archive/BPE_Gage.pdf}{\beamergotobutton{Gage (1994)}}}

	\begin{itemize}
		\item Considering data on a \textit{byte}-level
		\item Looking at pairs of bytes:
			\begin{enumerate}
				\item Count the occurrences of all byte pairs
				\item Find the most frequent byte pair
				\item Replace it with an unused byte
			\end{enumerate}
		\item Repeat this process until no further compression is possible
	\end{itemize}

\vfill

\end{vbframe}

% ------------------------------------------------------------------------------

\begin{vbframe}{BytePair encoding (BPE)}

\vfill

\textbf{Open-vocabulary neural machine translation \href{https://www.aclweb.org/anthology/P16-1162.pdf}{\beamergotobutton{Sennrich et al. (2016)}}}
	
	\begin{itemize}
		\item Instead of looking at bytes, look at characters
		\item Motivation: Translation as an open-vocabulary problem
		\item Word-level NMT models:
			\begin{itemize}
				\item Handling out-of-vocabulary word by using back-off dictionaries
				\item Unable to translate or generate previously unseen words
			\end{itemize}
		\item Using BPE effectively \textit{solves} this problem, except for ..
			\begin{itemize}
				\item .. the occurence of unknown characters
				\item .. when all occurences in the training set were merged into "larger" symbols
							(Example: \textit{"safeguar"} and \textit{"safeguard"})
			\end{itemize}
	\end{itemize}

\vfill

\end{vbframe}

% ------------------------------------------------------------------------------

\begin{vbframe}{BytePair encoding (BPE)}

\vfill

\textbf{Adapt BPE for word segmentation \href{https://www.aclweb.org/anthology/P16-1162.pdf}{\beamergotobutton{Sennrich et al. (2016)}}}

	\begin{itemize}
		\item \textit{Goal:} Represent an open vocabulary by a vocabulary of fixed size\\
					$\rightarrow$ Use variable-length character sequences 
		\item Looking at pairs of characters:
			\begin{enumerate}
				\item Initialize the the vocabulary with all characters plus end-of-word token
				\item Count occurrences and find the most frequent character pair,\\
							e.g. "A" and "B" (\warning Word boundaries are \textbf{not} crossed) \\
							\textit{[Side effect: Can be run on a dictionary w/ frequency counts]}
				\item Replace it with the new token "AB"
			\end{enumerate}
		\item Only one hyperparameter: Vocabulary size\\
					(Initial vocabulary + Specified no. of merge operations)\\
					$\rightarrow$ Repeat this process until given $|V|$ is reached
	\end{itemize}

\vfill

\end{vbframe}

% ------------------------------------------------------------------------------

\begin{vbframe}{Example -- Setup \href{https://www.aclweb.org/anthology/P16-1162.pdf}{\beamergotobutton{Sennrich et al. (2016)}}}

\begin{lstlisting}[language=Python]
import re, collections

def get_stats(vocab):
  pairs = collections.defaultdict(int)
  for word, freq in vocab.items():
    symbols = word.split()
    for i in range(len(symbols)-1):
      pairs[symbols[i],symbols[i+1]] += freq
  return pairs
	
def merge_vocab(pair, v_in):
  v_out = {}
  bigram = re.escape(' '.join(pair))
  p = re.compile(r'(?<!\S)' + bigram + r'(?!\S)')
  for word in v_in:
    w_out = p.sub(''.join(pair), word)
    v_out[w_out] = v_in[word]
  return v_out
\end{lstlisting}

\end{vbframe}

% ------------------------------------------------------------------------------

\begin{vbframe}{Example -- Merging \href{https://www.aclweb.org/anthology/P16-1162.pdf}{\beamergotobutton{Sennrich et al. (2016)}}}

\vspace{-.5cm}

\begin{lstlisting}[language=Python]
vocab = {'l o w </w>' : 5, 'l o w e r </w>' : 2, 
				 'n e w e s t </w>':6, 'w i d e s t </w>':3}

pairs = get_stats(vocab)
\end{lstlisting}

\vspace{-.75cm}

\begin{lstlisting}[language=Python]
>>> print(pairs)
defaultdict(<class 'int'>, {
						('l', 'o'): 7, ('o', 'w'): 7, ('w', '</w>'): 5, 
						('w', 'e'): 8, ('e', 'r'): 2, ('r', '</w>'): 2, 
						('n', 'e'): 6, ('e', 'w'): 6, ('e', 's'): 9, 
						('s', 't'): 9, ('t', '</w>'): 9, ('w', 'i'): 3, 
						('i', 'd'): 3, ('d', 'e'): 3
})
\end{lstlisting}

\vspace{-.5cm}

\begin{lstlisting}[language=Python]
best = max(pairs, key=pairs.get)
vocab = merge_vocab(best, vocab)
\end{lstlisting}

\vspace{-.75cm}

\begin{lstlisting}[language=Python]
>>> print(best)
('e', 's')
>>> print(vocab)
{'l o w </w>': 5, 'l o w e r </w>': 2, 
 'n e w es t </w>': 6, 'w i d es t </w>': 3}
\end{lstlisting}

\end{vbframe}

% ------------------------------------------------------------------------------

\begin{vbframe}{Example -- Merging \href{https://www.aclweb.org/anthology/P16-1162.pdf}{\beamergotobutton{Sennrich et al. (2016)}}}

\vspace{-.5cm}

\begin{lstlisting}[language=Python]
vocab = {'l o w </w>' : 5, 'l o w e r </w>' : 2, 
				 'n e w e s t </w>':6, 'w i d e s t </w>':3}

num_merges = 10

for i in range(num_merges):
  pairs = get_stats(vocab)
  best = max(pairs, key=pairs.get)
  vocab = merge_vocab(best, vocab)
  print(best)
\end{lstlisting}

\vspace{-.5cm}

\begin{lstlisting}[language=Python]
('e', 's')
('es', 't')
('est', '</w>')
('l', 'o')
('lo', 'w')
('n', 'e')
('ne', 'w')
('new', 'est</w>')
('low', '</w>')
('w', 'i')
\end{lstlisting}

\end{vbframe}

% ------------------------------------------------------------------------------

\endlecture
\end{document}
